% Generated draft: requires mathematical and source review.
\documentclass[11pt]{article}
\usepackage[margin=25mm]{geometry}
\usepackage{fontspec}
\setmainfont{Latin Modern Roman}
\usepackage{amsmath,amssymb,mathtools,unicode-math}
\setmathfont{Latin Modern Math}
\usepackage{enumitem,needspace,xurl,hyperref}
\setlength{\parindent}{0pt}
\setlength{\parskip}{5pt}
\setlength{\emergencystretch}{3em}
\newcommand{\N}{\mathbb{N}}
\newcommand{\Z}{\mathbb{Z}}
\newcommand{\Q}{\mathbb{Q}}
\newcommand{\R}{\mathbb{R}}
\newcommand{\C}{\mathbb{C}}
\newcommand{\F}{\mathbb{F}}
\newcommand{\A}{\mathbb{A}}
\newcommand{\PP}{\mathbb{P}}
\newcommand{\E}{\mathbb{E}}
\newcommand{\eps}{\varepsilon}
\newcommand{\dd}{\,\mathrm{d}}
\begin{document}
\section*{UnsolvedMath ID: 20002815}
\textbf{Source problem code:} catalogue entry
\subsection{Definitions and mathematical statement}
The source asks whether Andr\'e's supercharacter theory is the coarsest retaining the single-box elementary irreducibles. Its reference to irreducible ``superclasses'' is imprecise, since irreducibility is a character property, and ``the irreducibles'' could mean this distinguished family or all irreducible characters. The aid uses the former reading and retains every field-labeled elementary character separately.

For a finite group $G$, a complex supercharacter theory consists of partitions $\mathcal X$ of $\operatorname{Irr}(G)$ and $\mathcal K$ of $G$ into superclasses, with the standard identity, equal-cardinality, and constancy conditions for
\[
 \sigma_X=\sum_{\chi\in X}\chi(1)\chi.
\]
Its superclass-function space $\operatorname{SCF}$ is the unital pointwise algebra of functions constant on $\mathcal K$, spanned by these supercharacters. To retain an irreducible $\chi$ means that $\{\chi\}$ is a character block, not merely that it occurs as a constituent somewhere. Coarsest means that every competing theory retaining the specified family refines its superclass and character partitions.

The aid's general criterion is: if a supercharacter theory $\mathcal S_0$ retains a family $E$ of singleton irreducibles and
\[
 \operatorname{SCF}(\mathcal S_0)=\C[E]
\]
as a unital pointwise algebra, then it is the unique coarsest theory retaining $E$. Its superclasses are exactly the joint value fibers
\[
 g\sim h\quad\Longleftrightarrow\quad
 \chi(g)=\chi(h)\text{ for every }\chi\in E.
\]
Indeed the comparison concerns the function algebra generated by fixed labeled characters, not unlabeled or phase-modified value vectors.

For $G=U_n(q)$, the upper-unitriangular group over $\F_q$, take $n\ge1$ and all elementary characters $\chi_{ij,a}$ with $1\le i<j\le n$ and $a\in\F_q^\times$. The aid reports that these are singleton irreducibles in the Andr\'e--Yan theory, and every basic supercharacter associated to a labeled rook placement is the pointwise product of its single-box elementary characters. Since the basic supercharacters span its superclass-function space, the criterion gives the asserted unique coarsest property for every prime power $q$. This statement uses all labels and all allowed single positions, not only first-superdiagonal linear characters.

If the final source phrase instead requires every irreducible character of $G$ to be a singleton block, the theory is simply ordinary character theory with conjugacy classes as superclasses. That is a different reading. The aid's reported solution addresses the specified elementary-family interpretation; the exact original intent and cited product theorem still require review.
\subsection{Short English statement}
Under the explicit reading that all labeled elementary irreducibles must remain singleton blocks, the aid reports that their pointwise products generate Andr\'e--Yan's entire superclass algebra, making it the unique coarsest such theory. Requiring all irreducibles instead would give ordinary character theory.
\subsection{Sources}
\begin{itemize}
\item Local UnsolvedMath registry, permanent ID 20002815, source code catalogue entry; created and updated 14 August 2026. Canonical \texttt{aim-representation-theory-notes.json}, zero-based \texttt{notes[86]}, tag \texttt{problem}, item 1.3.
\item AIM ``Supercharacters and combinatorial Hopf algebras'': \url{https://aimath.org/WWN/supercharacters/supercharacters.pdf}. Wording is marked exact; original remarks and literature field are empty. The source names Andr\'e's theory without a complete bibliographic citation.
\item Research attempt 1: \texttt{full\_solution}, problem status at run solved, rigor label \texttt{proved}, novelty confidence low. Registry labels are solved and SOLVED-IN-LITERATURE. No complete authoritative reference or product-formula theorem statement for the Andr\'e--Yan theory is supplied.
\item Ulam AI / UnsolvedMath Contributors, CC BY 4.0. Machine-generated research aid, not peer reviewed; independently verify before citation.
\end{itemize}
\end{document}
